\documentclass[11pt,a4paper]{article}
\usepackage[margin=0.75in]{geometry}
\usepackage{amsmath,amssymb,mathtools,amsthm}
\usepackage{booktabs}
\usepackage{tabularx}
\usepackage{longtable}
\usepackage{array}
\usepackage{enumitem}
\usepackage{hyperref}
\usepackage{graphicx}
\usepackage{xcolor}
\usepackage{caption}
\usepackage{float}
\usepackage{fancyvrb}
\usepackage{tikz}
\usetikzlibrary{arrows.meta,positioning,shapes.geometric,calc}
\usepackage{ctex} % 中文支持

% ===== 宏定义 =====
\newcommand{\RR}{\mathbb{R}}
\newcommand{\sigmoid}{\sigma}
\newcommand{\clip}{\operatorname{clip}}
\newcommand{\normop}{\operatorname{norm}}
\newcommand{\IoU}{\operatorname{IoU}}
\newcommand{\cosine}{\operatorname{cos}}

\theoremstyle{definition}
\newtheorem{definition}{定义}[section]
\newtheorem{proposition}{命题}[section]
\newtheorem{remark}{备注}[section]

% >>> 修改：移除标题末尾多余的 "ACM" 文字
\title{\textbf{ProveTok v3.0：完整数学工作流}\\[6pt]
\large 方法 M1--M5 全推导、核验与流水线总结}
\date{\today}

\begin{document}
\maketitle
\tableofcontents
\newpage

%=================================================================
\section{执行摘要}
%=================================================================

ProveTok v3.0 是一个六阶段流水线，用于生成带句级证据追溯的可验证 3D CT 放射报告。该流水线将 CT 体数据转化为固定大小的\emph{证据 token bank}，通过轻量跨模态路由器将 token 子集路由到每句报告，在 token 门控接口下生成句子（构造性地保证引用正确性），应用规则校验器，并在统计稳健的评估协议下报告结果。

本文档提供每个阶段（方法 M1--M5）的\textbf{完整数学规范}，包括逐步推导、边界情形分析、单调性/有界性证明，以及统一的符号参考。

%=================================================================
\section{统一符号表}
%=================================================================

\begin{table}[ht]
\centering\small
\begin{tabularx}{\linewidth}{lX}
\toprule
\textbf{符号} & \textbf{描述} \\
\midrule
$V \in \RR^{D \times H \times W}$ & 输入 CT 体数据 \\
$F = \{f_i\}_{i=1}^{N}$ & 冻结 3D 编码器输出的特征图；$f_i \in \RR^d$ \\
$A_i \in [0,1]$ & 第 $i$ 个空间单元的伪影风险分数 \\
$g_i \in (0,1)$ & 由 $A_i$ 导出的 sigmoid 软门控 \\
$H_i$ & 第 $i$ 个单元的不确定性（熵） \\
$P_i$ & 第 $i$ 个单元的语义先验分数 \\
$q(\cdot) \in [0,1]$ & 分位数归一化（按深度层、按病例） \\
$S_i \in [0,1)$ & 第 $i$ 个单元的综合重要性分数 \\
$B$ & token bank 大小（默认 64） \\
$k$ & 每句 top-$k$ 证据 token（默认 8） \\
$W_{\mathrm{proj}} \in \RR^{d_t \times d}$ & 可学习线性投影矩阵（路由器） \\
$v_i \in \RR^{d_t}$ & 投影并归一化后的视觉 token 嵌入 \\
$q_s \in \RR^{d_t}$ & 句子主题 $s$ 的文本查询嵌入 \\
$r_i^{(s)}$ & token $i$ 对句子 $s$ 的余弦路由分数 \\
$\hat{r}_i^{(s)}$ & 加入空间先验后的路由分数 \\
$\tau$ & InfoNCE 温度参数（默认 0.07） \\
$\mathcal{P}_s, \mathcal{N}_s$ & 句子 $s$ 的正样本 / 负样本 token 集合 \\
$\mathrm{sev}_i$ & 规则违反的严重度权重 \\
% >>> 新增符号
$\gamma$ & 对数平滑惩罚缩放因子（默认 2.0） \\
$\epsilon$ & 解剖优先模式中语义 tiebreak 权重（默认 0.05） \\
$\theta_{\mathrm{ratio}}$ & R1 侧别同侧 ratio 阈值（默认 0.6） \\
$\theta_{\mathrm{support}}$ & R2 解剖支持率阈值（默认 0.8） \\
\bottomrule
\end{tabularx}
\caption{ProveTok v3.0 统一符号表。}
\label{tab:notation}
\end{table}


%=================================================================
\section{流水线总览}
%=================================================================

\begin{figure}[ht]
\centering
\begin{tikzpicture}[
    node distance=0.9cm and 1.2cm,
    block/.style={rectangle, draw, rounded corners, fill=blue!8,
        minimum width=3.4cm, minimum height=0.85cm, align=center, font=\small},
    decision/.style={diamond, draw, fill=yellow!15, minimum width=1.6cm,
        minimum height=0.8cm, align=center, font=\small, inner sep=1pt},
    arrow/.style={-{Stealth[length=2.5mm]}, thick},
]
\node[block] (vol) {CT 体数据 $V$};
\node[block, below=of vol] (s0) {Stage 0: 伪影风险 $A_i$};
\node[block, right=2.5cm of s0] (s1) {Stage 1: 冻结 3D 编码器 $\to F$};
\node[block, below=of s0, xshift=2cm] (s2) {Stage 2: 自适应分裂 $\to$ $B$ 个 tokens};
\node[block, below=of s2] (s3a) {Stage 3a: 证据规划};
\node[block, below=of s3a] (s3b) {Stage 3b: 路由器 top-$k$};
\node[block, below=of s3b] (s3c) {Stage 3c: token 门控生成};
\node[decision, below=1.0cm of s3c] (s4) {Stage 4:\\校验器};
\node[block, right=2.0cm of s4] (s5) {Stage 5: 重路由};
\node[block, below=1.0cm of s4] (s6) {Stage 6: 输出 + 追溯日志};
\draw[arrow] (vol) -- (s0);
\draw[arrow] (vol) -| (s1);
\draw[arrow] (s0) -- (s2);
\draw[arrow] (s1) |- (s2);
\draw[arrow] (s2) -- (s3a);
\draw[arrow] (s3a) -- (s3b);
\draw[arrow] (s3b) -- (s3c);
\draw[arrow] (s3c) -- (s4);
\draw[arrow] (s4) -- node[left, font=\small]{通过} (s6);
\draw[arrow] (s4) -- node[above, font=\small]{失败} (s5);
\draw[arrow] (s5) |- (s3b);
\end{tikzpicture}
\caption{ProveTok v3.0 六阶段流水线。Stage 0--2 = \textbf{M1}（token bank），Stage 3b = \textbf{M2}（路由器），Stage 3c = \textbf{M3}（门控生成），Stage 4--5 = \textbf{M4}（校验+重路由），评估协议 = \textbf{M5}。}
\label{fig:pipeline}
\end{figure}


%=================================================================
\section{M1：3D 证据 Token Bank 构建（Stage 0--2）}
%=================================================================

\subsection{Stage 0：伪影风险评分}

\begin{definition}[伪影风险分数]
对每个空间单元 $i$：
\begin{equation}\label{eq:artifact}
A_i = \clip\!\Bigg(
    w_{\mathrm{snr}} \cdot \normop\!\!\left(\frac{\sigma_i}{|\mu_i|+\epsilon}\right)
    + w_{\mathrm{st}} \cdot \normop(\mathrm{streak}_i)
    + w_{\mathrm{out}} \cdot \normop(\mathrm{outlier}_i),\;0,\;1
\Bigg),
\end{equation}
其中 $\mu_i, \sigma_i$ 为局部均值与标准差，$\mathrm{streak}_i$ 捕捉条纹状结构噪声，$\mathrm{outlier}_i$ 衡量体素级异常程度，$\normop(\cdot)$ 将各指标映射到 $[0,1]$。权重满足 $w_{\mathrm{snr}} + w_{\mathrm{st}} + w_{\mathrm{out}} = 1$，默认 $(0.5, 0.3, 0.2)$。
\end{definition}

\paragraph{输入 / 输出 / 参数来源。}
\begin{itemize}[nosep]
    \item \textbf{输入：}CT volume $V$ 中第 $i$ 个空间单元（八叉树叶节点）的局部体素统计。$\mu_i, \sigma_i$ 由该单元对应的体素灰度值计算；$\mathrm{streak}_i$ 通过方向梯度检测条纹结构；$\mathrm{outlier}_i$ 通过体素灰度 z-score 阈值判定。
    \item \textbf{输出：}$A_i \in [0,1]$ — 伪影风险概率。传给下游 sigmoid 软门控（公式~\ref{eq:gate}）。
    \item \textbf{参数来源：}权重 $(w_{\mathrm{snr}}, w_{\mathrm{st}}, w_{\mathrm{out}}) = (0.5, 0.3, 0.2)$ 为预设超参，根据 CT 伪影的临床先验设定（SNR 异常最具区分力，条纹次之）。$\epsilon$ 为数值稳定常数（防止除零）。
\end{itemize}

\paragraph{设计动机。}
逆信噪比项 $\sigma_i/(|\mu_i|+\epsilon)$ 在噪声主导时增大。条纹伪影（金属/射束硬化）与统计离群值提供互补检测通道。裁剪到 $[0,1]$ 保证 $A_i$ 为有效的风险概率。

% >>> 新增：明确 norm(·) 实现
\begin{remark}[$\normop(\cdot)$ 实现]
代码中 $\normop(\cdot)$ 采用 min-max 归一化：对同一批次的原始值 $\{x_i\}$，$\normop(x_i) = (x_i - x_{\min}) / (x_{\max} - x_{\min})$。当 $x_{\max} = x_{\min}$ 时回退到常数 $0.5$。
\end{remark}


\subsection{Sigmoid 软门控}

\begin{definition}[软门控]
\begin{equation}\label{eq:gate}
g_i = \sigmoid\!\big(-k_{\mathrm{gate}}(A_i - \tau_A)\big)
    = \frac{1}{1 + \exp\!\big(k_{\mathrm{gate}}(A_i - \tau_A)\big)},
\end{equation}
其中 $k_{\mathrm{gate}} > 0$ 控制陡峭度（默认 15），$\tau_A \in (0,1)$ 为伪影阈值（默认 0.85）。
\end{definition}

\paragraph{输入 / 输出 / 参数来源。}
\begin{itemize}[nosep]
    \item \textbf{输入：}$A_i \in [0,1]$（来自公式~\ref{eq:artifact} 的伪影风险分数）。
    \item \textbf{输出：}$g_i \in (0,1)$ — 软门控权重。$g_i \to 1$ 表示``低伪影，保留''；$g_i \to 0$ 表示``高伪影，抑制''。传给下游综合评分 $S_i$（公式~\ref{eq:score}）。
    \item \textbf{参数来源：}$\tau_A = 0.85$ 为伪影阈值，基于临床 CT 图像质量标准（$A_i > 0.85$ 的区域在视觉上不可接受）。$k_{\mathrm{gate}} = 15$ 控制过渡陡峭度——取 15 使阈值两侧 $\pm 0.1$ 范围内的门控值从 0.82 迅速跌至 0.18，实现近似硬截断同时保持可微。
\end{itemize}

\begin{proposition}[门控性质]\label{prop:gate}
\begin{enumerate}[label=(\roman*)]
    \item \textbf{值域：}对所有 $A_i \in [0,1]$，$g_i \in (0,1)$。
    \item \textbf{单调性：}$g_i$ 关于 $A_i$ 严格递减。
    \item \textbf{阈值行为：}$g_i(\tau_A) = 0.5$；$A_i \gg \tau_A \Rightarrow g_i \to 0$；$A_i \ll \tau_A \Rightarrow g_i \to 1$。
\end{enumerate}
\end{proposition}

\begin{proof}
(i) 由于 $\sigmoid(x) \in (0,1)$ 对所有 $x \in \RR$ 成立，值域立即得证。

(ii) 对 $A_i$ 求导：
\begin{equation}
\frac{\partial g_i}{\partial A_i}
= \underbrace{\sigmoid(-k_{\mathrm{gate}}(A_i-\tau_A))}_{>0}\;
  \underbrace{(1 - \sigmoid(-k_{\mathrm{gate}}(A_i-\tau_A)))}_{>0}
  \cdot \underbrace{(-k_{\mathrm{gate}})}_{<0} < 0.
\end{equation}
因此 $g_i$ 关于 $A_i$ 严格递减。

(iii) 当 $A_i=\tau_A$ 时：$g_i=\sigmoid(0)=0.5$。取默认值 $k_{\mathrm{gate}}=15,\,\tau_A=0.85$：
\[
A_i=1 \Rightarrow g_i \approx \sigmoid(-2.25) \approx 0.095, \quad
A_i=0 \Rightarrow g_i \approx \sigmoid(12.75) \approx 1.0.
\]
\end{proof}


\subsection{Stage 1：冻结 3D 编码器}

预训练的 3D 编码器（SwinUNETR，$\mathrm{feature\_size}=48$，$\mathrm{spatial\_dims}=3$）处理体数据 $V$，输出特征图 $F = \{f_i \in \RR^d\}_{i=1}^{N}$。输入体积统一 resize 到 $128^3$。编码器权重\textbf{全程冻结}，无梯度回传。

\paragraph{输入 / 输出 / 参数来源。}
\begin{itemize}[nosep]
    \item \textbf{输入：}CT volume $V$，经预处理 resize 至 $128 \times 128 \times 128$ 体素。
    \item \textbf{输出：}特征图 $F = \{f_i \in \RR^{768}\}_{i=1}^{N}$，其中 $d=768$ 由 SwinUNETR 的 \texttt{feature\_size=48} 决定（bottleneck 输出维度 $= 48 \times 16 = 768$）。$N$ 为空间网格大小（初始 $128^3 / \text{patch}^3$）。
    \item \textbf{参数来源：}SwinUNETR 权重来自 BTCV（Beyond the Cranial Vault）多器官分割预训练检查点（MONAI 提供），全程冻结以避免在小数据集上过拟合并保证特征一致性。
\end{itemize}


\subsection{Stage 2：综合评分}

\begin{definition}[分位数归一化]
对同一八叉树深度层、同一病例内的原始值 $\{x_i\}$：
\begin{equation}
q(x_i) = \frac{\mathrm{rank}(x_i)}{n}, \quad n = |\{x_j : \text{同一深度层}\}|.
\end{equation}
在每轮处理中冻结（迭代分裂时不重算）。对离群值鲁棒，输出近似均匀分布。
\end{definition}

\begin{definition}[综合重要性分数]
\begin{equation}\label{eq:score}
S_i = g_i \cdot \big(\lambda_h \, q(H_i) + \lambda_p \, q(P_i)\big), \quad \lambda_h + \lambda_p = 1.
\end{equation}
默认：$\lambda_h = 0.6$，$\lambda_p = 0.4$。

\paragraph{输入 / 输出 / 参数来源。}
\begin{itemize}[nosep]
    \item \textbf{输入：}$g_i$（来自公式~\ref{eq:gate}）、$H_i$（不确定性，包括特征熵 $H_i^{\mathrm{ent}}$、patch 方差 $H_i^{\mathrm{var}}$、边界指示 $B_i$）、$P_i$（语义先验，由 $F$ 的局部范数 $\|F\|_{\Omega_i}$ 经 sigmoid 门控导出）。均在 Stage 1 输出 $F$ 上计算。
    \item \textbf{输出：}$S_i \in [0,1)$ — 综合重要性分数。用于自适应八叉树分裂的候选排序和 top-$B$ 选择。
    \item \textbf{参数来源：}$\lambda_h = 0.6, \lambda_p = 0.4$ — 不确定性优先于语义先验，因为高不确定性区域更可能包含需要详细报告的病理性发现。$q(\cdot)$ 按同一深度层的分位数归一化，确保不同深度层的分数可比。
\end{itemize}
\end{definition}

\begin{proposition}[分数有界性]\label{prop:score}
对所有合法输入，$S_i \in [0, 1)$。
\end{proposition}
\begin{proof}
由 $q(\cdot) \in [0,1]$ 及 $\lambda_h + \lambda_p = 1$：
\[
0 \leq \lambda_h q(H_i) + \lambda_p q(P_i) \leq 1.
\]
由 $g_i \in (0,1)$（命题~\ref{prop:gate}）：
\[
0 \leq S_i = g_i \cdot (\lambda_h q(H_i) + \lambda_p q(P_i)) < 1.
\]
下界在 $q(H_i)=q(P_i)=0$ 时取到。上界严格小于 1，因为 $g_i < 1$。
\end{proof}


\subsection{自适应八叉树分裂算法}

\begin{enumerate}
    \item \textbf{初始化：}以整个体数据为单个根单元。
    \item \textbf{选择候选：}在当前叶子中选 $i^* = \arg\max_i S_i$（需超过分裂阈值）。
    \item \textbf{分裂：}将 $i^*$ 细分为 8 个子单元；对每个子 $j$ 计算 $f_j, A_j, g_j, H_j, P_j, S_j$。
    \item \textbf{上限检查：}若叶子数 $\geq B$，执行 batched top-$B$ + 空间 NMS 裁剪到恰好 $B$ 个 token。
    \item \textbf{终止：}无候选超过阈值，或已达上限 $B$。
\end{enumerate}

该过程类似于 OctNet 的不平衡八叉树（叶节点存 pooled feature），但以 $S_i$（而非 occupancy）作为分裂准则，将空间分辨率集中在高重要性区域。


\subsection{边界感知导出}

\begin{definition}[边界上下文混合]
对有邻居 $\mathcal{N}(i)$ 的导出 token $i$：
\begin{align}
b_i &= \frac{1}{|\mathcal{N}(i)|} \sum_{j \in \mathcal{N}(i)} f_j, \label{eq:neighbor}\\
f_i^{\mathrm{export}} &= (1 - \beta)\, f_i + \beta\, b_i, \quad \beta \in [0,1]. \label{eq:export}
\end{align}
默认 $\beta = 0.1$。边界框额外扩展 12.5\%。

\paragraph{输入 / 输出 / 参数来源。}
\begin{itemize}[nosep]
    \item \textbf{输入：}八叉树选中的 $B$ 个叶节点特征 $\{f_i\}$，以及每个节点的邻居集合 $\mathcal{N}(i)$（同深度、空间相邻）。
    \item \textbf{输出：}\textbf{Token Bank} $\{(f_i^{\mathrm{export}}, \mathrm{bbox}_i, \mathrm{level}_i)\}_{i=1}^{B}$ — 即 $B$ 个带空间坐标和深度层级的视觉 token。这是整个 M1 模块的最终输出，直接传给 M2 路由器。
    \item \textbf{参数来源：}$\beta = 0.1$ 为轻度邻域融合，仅在边界处引入少量空间上下文以减少 block artifact。$12.5\%$ bbox 扩展为工程选择，确保 token 边界略有重叠。
\end{itemize}
\end{definition}

\begin{proposition}[凸组合稳定性]
\begin{equation}
\|f_i^{\mathrm{export}}\| \leq (1-\beta)\|f_i\| + \beta\|b_i\|
\leq (1-\beta)\|f_i\| + \frac{\beta}{|\mathcal{N}(i)|}\sum_{j \in \mathcal{N}(i)} \|f_j\|.
\end{equation}
若所有特征范数以 $M$ 为界，则 $\|f_i^{\mathrm{export}}\| \leq M$。不会数值爆炸。
\end{proposition}

\begin{remark}[邻域为空的回退]
当 $|\mathcal{N}(i)|=0$（孤立叶节点）：令 $b_i = f_i$（等价 $\beta=0$），得 $f_i^{\mathrm{export}}=f_i$。
\end{remark}

\begin{remark}[归一化交互]
若 $f_i$ 经 $L_2$ 归一化，线性均值 $b_i$ 可能改变方向与范数。建议：(a) 在归一化前混合，或 (b) 混合后重新归一化。需在论文中明确说明选择。
\end{remark}


%=================================================================
\section{M2：轻量跨模态路由器}
%=================================================================

\subsection{线性投影与余弦路由}

\begin{definition}[视觉 token 投影]
\begin{equation}\label{eq:proj}
v_i = \frac{W_{\mathrm{proj}} f_i}{\|W_{\mathrm{proj}} f_i\|_2}, \quad v_i \in \RR^{d_t}, \quad \|v_i\|_2 = 1.
\end{equation}
仅 $W_{\mathrm{proj}} \in \RR^{d_t \times d}$ 可学习。

\paragraph{输入 / 输出 / 参数来源。}
\begin{itemize}[nosep]
    \item \textbf{输入：}$f_i \in \RR^{768}$（来自 M1 Token Bank 的视觉特征 $f_i^{\mathrm{export}}$）。
    \item \textbf{输出：}$v_i \in \RR^{384}$，$\|v_i\|_2 = 1$ — 归一化的视觉嵌入，与文本查询在同一 $384$ 维空间中。
    \item \textbf{参数来源：}$W_{\mathrm{proj}} \in \RR^{384 \times 768}$ 是\textbf{整个流水线中唯一可学习的参数}（$384 \times 768 = 294{,}912$ 个参数）。$d_t = 384$ 由文本编码器 \texttt{all-MiniLM-L6-v2} 的输出维度决定（必须匹配）。初始化为正交矩阵。当未训练时退化为截断单位矩阵（A0 配置）。
\end{itemize}
\end{definition}

\begin{definition}[路由分数]
给定 $L_2$ 归一化的查询 $q_s$（$\|q_s\|_2=1$）：
\begin{equation}\label{eq:route}
r_i^{(s)} = \cosine(q_s, v_i) = q_s^\top v_i \in [-1, 1].
\end{equation}

\paragraph{输入 / 输出 / 参数来源。}
\begin{itemize}[nosep]
    \item \textbf{输入：}$q_s \in \RR^{384}$（来自 \texttt{all-MiniLM-L6-v2} 冻结文本编码器，输入为句子级 topic 字符串，如 ``consolidation in right lower lobe''）和 $v_i \in \RR^{384}$（来自公式~\ref{eq:proj}）。两者均已 $L_2$ 归一化。
    \item \textbf{输出：}$r_i^{(s)} \in [-1, 1]$ — 每个 token $i$ 对句子 $s$ 的语义相关度。值越大表示该 token 的视觉特征与句子主题在嵌入空间中越对齐。
    \item \textbf{参数来源：}无额外参数。余弦相似度是标准度量；$L_2$ 归一化保证值域在 $[-1,1]$。
\end{itemize}
\end{definition}

\begin{definition}[空间先验增强]
当解剖关键词的边界框可用时：
\begin{equation}\label{eq:spatial}
\hat{r}_i^{(s)} = r_i^{(s)} + \lambda_{\mathrm{spatial}} \cdot \IoU(\mathrm{bbox}_i, \mathrm{bbox}_{\mathrm{anatomy}}^{(s)}),
\end{equation}
其中 $\IoU \in [0,1]$，$\lambda_{\mathrm{spatial}} \geq 0$（默认 0.3）。因此 $\hat{r}_i^{(s)} \in [-1, 1+\lambda_{\mathrm{spatial}}]$。

\paragraph{输入 / 输出 / 参数来源。}
\begin{itemize}[nosep]
    \item \textbf{输入：}$r_i^{(s)}$（公式~\ref{eq:route}）+ $\mathrm{bbox}_i$（token 的 3D 边界框，来自 M1 Token Bank）+ $\mathrm{bbox}_{\mathrm{anatomy}}^{(s)}$（句子 $s$ 对应解剖区域的边界框，由 Stage 3a 证据规划器输出的解剖关键词查 \texttt{DEFAULT\_ANATOMY\_BOXES} 规则表获得，坐标为归一化 $[0,1]$ 后乘以 volume 维度 $128^3$）。
    \item \textbf{输出：}$\hat{r}_i^{(s)}$ — 融合空间先验后的路由分数。对所有 $B$ 个 token 排序后取 top-$k$（默认 $k=8$），得到证据集 $E_s$。
    \item \textbf{参数来源：}$\lambda_{\mathrm{spatial}} = 0.3$ 为我们实验中设定的超参，通过消融网格 $\{0, 0.1, 0.3, 0.5\}$ 选出（见表~\ref{tab:hyper}）。在解剖优先模式~\eqref{eq:anat_primary} 中，$\epsilon = 0.05$ 确保同 IoU token 之间仍有语义分辨力。
\end{itemize}
\end{definition}

按 $\hat{r}_i^{(s)}$ 取 top-$k$ 的 token 构成句子 $s$ 的证据集 $E_s$。

% >>> 新增：解剖优先路由模式
\begin{remark}[解剖优先路由模式]\label{rmk:anat_primary}
当 $W_{\mathrm{proj}}$ 尚未经过 InfoNCE 训练时（退化为单位矩阵），语义余弦分数 $r_i^{(s)}$ 不可靠。此时采用\textbf{解剖优先}模式：
\begin{equation}\label{eq:anat_primary}
\hat{r}_i^{(s)} = \IoU(\mathrm{bbox}_i,\, \mathrm{bbox}_{\mathrm{anatomy}}^{(s)}) + \epsilon \cdot r_i^{(s)},
\end{equation}
其中 $\epsilon \ll 1$（默认 $0.05$）使语义分数仅作为同 IoU token 间的 tiebreaker。该模式将空间 IoU 作为唯一有校准意义的信号，在 $W_{\mathrm{proj}}$ 训练完成后可切回标准模式~\eqref{eq:spatial}。
\end{remark}


\subsection{InfoNCE 训练损失：完整推导}

\begin{definition}[InfoNCE 损失]
令 $\mathcal{C}_s = \mathcal{P}_s \cup \mathcal{N}_s$：
\begin{equation}\label{eq:infonce}
\mathcal{L}_s = -\frac{1}{|\mathcal{P}_s|} \sum_{i \in \mathcal{P}_s}
\log \frac{\exp(\hat{r}_i^{(s)}/\tau)}
     {\sum_{j \in \mathcal{C}_s} \exp(\hat{r}_j^{(s)}/\tau)}.
\end{equation}
默认 $\tau = 0.07$。
\end{definition}

\paragraph{正样本与负样本的构建。}

训练数据通过 \textbf{bootstrap 策略}构建。具体步骤如下：

\subparagraph{Step 1：用 A0 配置生成伪标签。}
令 $W_{\mathrm{proj}} = I$（单位矩阵），跑一遍完整流水线（A0 配置）。对每个病例的每个句子 $s$：
\begin{enumerate}[nosep]
    \item Stage 3a 根据句子的 topic 文本查询解剖关键词表（\texttt{DEFAULT\_ANATOMY\_BOXES}），得到该句子对应的解剖区域 3D 边界框 $\mathrm{bbox}_{\mathrm{anatomy}}^{(s)}$。例如 topic 为 ``consolidation in right lower lobe'' 时，查到 right lower lobe 的标准 bbox。
    \item 将 $\mathrm{bbox}_{\mathrm{anatomy}}^{(s)}$ 与 Token Bank 中所有 $B=128$ 个 token 的 bbox 逐一计算 3D IoU。IoU 高的 token 在空间上落在该解剖区域内。
    \item 按 IoU 分数排序，取前 $k=8$ 个 token，将其 ID 写入 \texttt{trace.jsonl} 的 \texttt{topk\_token\_ids} 字段。
\end{enumerate}
由于 $W_{\mathrm{proj}} = I$ 时语义分数 $r_i^{(s)}$ 基本是噪声，排序实际由 IoU 项主导。因此正样本的``标签''不来自人工标注，而来自空间 IoU 的自动对齐——这就是 bootstrap 的含义：用一个弱信号（空间重叠）生成伪标签来训练一个更强的信号（语义匹配）。

\subparagraph{Step 2：从 trace.jsonl 提取正样本。}
InfoNCE 训练脚本读取 A0 产出的 \texttt{trace.jsonl}，对每个句子 $s$，将其 \texttt{topk\_token\_ids} 中记录的 $k$ 个 token 作为正样本集：
\[
\mathcal{P}_s = \{t_1, t_2, \ldots, t_k\}, \quad |\mathcal{P}_s| = k = 8.
\]
这 $k$ 个 token 就是与句子 $s$ 所描述的解剖区域空间重叠最大的视觉 token。

\subparagraph{Step 3：In-batch Negatives 构建负样本。}
一个 mini-batch 从训练集中采样 $N=32$ 个句子 $\{s_1, s_2, \ldots, s_N\}$，每个句子各有其正样本集。对句子 $s$ 而言，batch 内其余 $N-1$ 个句子的正样本 token 均作为负样本：
\[
\mathcal{N}_s = \bigcup_{s' \neq s} \mathcal{P}_{s'}, \quad
|\mathcal{N}_s| = (N-1) \times k = 31 \times 8 = 248.
\]
因此候选集大小为 $|\mathcal{C}_s| = |\mathcal{P}_s| + |\mathcal{N}_s| = k + (N-1)k = Nk = 256$。

\textbf{为什么用 in-batch negatives？}(1) 计算高效，无需额外采样，batch 内自然提供；(2) 同一 batch 的句子来自不同解剖区域，形成``硬负样本''——例如 mediastinum 的 token 和 left lung 的 token 在空间上可能很近，迫使模型学到语义区分而非单纯空间距离；(3) 与 MoCo/SimCLR/CLIP 的训练范式一致。

\textbf{False negative 问题。}若 batch 内两个句子恰好描述同一解剖区域（如 $s_1$ 和 $s_5$ 都描述 left lung），则 $s_5$ 的正样本 token 被错误地当作 $s_1$ 的负样本。但实践中 batch 内 32 个句子通常来自不同病例的不同解剖区域，碰撞概率很低，且 InfoNCE 对少量 false negatives 具有鲁棒性~(Chuang et al., 2020)。

\paragraph{训练目标的直觉。}
InfoNCE 训练使 $W_{\mathrm{proj}}$ 学会将句子查询 $q_s$ 与其正样本 token（空间上对应的 token）在 $\RR^{384}$ 嵌入空间中拉近，同时推远与其他句子关联的 token（负样本）。训练完成后，语义余弦分数 $r_i^{(s)}$ 变得可靠，可以从解剖优先模式（公式~\ref{eq:anat_priority}）切换回标准模式（公式~\ref{eq:spatial}）。

\paragraph{输入 / 输出 / 参数来源。}
\begin{itemize}[nosep]
    \item \textbf{输入：}A0 流水线产出的 \texttt{trace.jsonl}（提供 $(s, \mathcal{P}_s)$ 配对），以及 M1 Token Bank 导出的冻结视觉特征 $\{f_i^{\mathrm{export}} \in \RR^{768}\}$。
    \item \textbf{输出：}训练后的 $W_{\mathrm{proj}} \in \RR^{384 \times 768}$ — 使视觉嵌入 $v_i$ 与语义相关的文本查询 $q_s$ 在 $\RR^{384}$ 空间中对齐。
    \item \textbf{参数来源：}$\tau = 0.07$ 来自 MoCo/CLIP 的标准温度值；batch size $N=32$；AdamW 优化器（$\mathrm{lr}=10^{-3}$，$\mathrm{weight\_decay}=10^{-4}$）；余弦退火调度器；早停 patience=10 epochs（最大 50 epochs）。
\end{itemize}

\subsubsection{训练过程：前向传播}

给定 batch 内句子 $s$ 和候选集 $\mathcal{C}_s$ 中的所有 token $j$，前向传播分为以下步骤：

\paragraph{Step 3a — Token 投影。}
每个 token 的冻结视觉特征 $f_j \in \RR^{768}$（来自 M1 Token Bank，不参与梯度更新）经过可学习投影矩阵并 $L_2$ 归一化：
\[
f_j \in \RR^{768} \xrightarrow{W_{\mathrm{proj}}} W_{\mathrm{proj}} f_j
\xrightarrow{L_2\text{ norm}} v_j = \frac{W_{\mathrm{proj}} f_j}{\|W_{\mathrm{proj}} f_j\|_2} \in \RR^{384}.
\]
$W_{\mathrm{proj}} \in \RR^{384 \times 768}$ 是整个训练过程中\textbf{唯一可学习的参数}（$384 \times 768 = 294{,}912$ 个参数）。

\paragraph{Step 3b — 文本查询编码。}
\[
q_s = \mathrm{MiniLM}(\mathrm{topic}_s) \in \RR^{384}, \quad \|q_s\|_2 = 1.
\]
MiniLM 编码器同样冻结，不参与梯度更新。

\paragraph{Step 3c — 计算余弦相似度（路由分数）。}
\[
r_j^{(s)} = q_s^\top v_j \in [-1, 1].
\]

\paragraph{Step 3d — 温度缩放 + Softmax。}
\begin{equation}
z_j \triangleq \frac{r_j^{(s)}}{\tau}, \qquad
p_j \triangleq \frac{\exp(z_j)}{\sum_{u \in \mathcal{C}_s} \exp(z_u)}.
\end{equation}
满足 $\sum_j p_j = 1$，$p_j > 0$。$\tau = 0.07$ 很小，使得 softmax 分布非常尖锐：两个 token 余弦差 0.1 时，logit 差为 $0.1/0.07 \approx 1.43$，softmax 概率比差 $e^{1.43} \approx 4.2$ 倍。这迫使模型精确区分正负样本。

\paragraph{Step 3e — 计算 InfoNCE 损失。}
对每个正样本 $i \in \mathcal{P}_s$：
\begin{equation}
\ell_i = -\log p_i = -\frac{r_i^{(s)}}{\tau} + \log \sum_{u \in \mathcal{C}_s} \exp\!\Big(\frac{r_u^{(s)}}{\tau}\Big).
\end{equation}
对句子 $s$ 的总损失（多正样本取平均）：
\begin{equation}
\mathcal{L}_s = \frac{1}{|\mathcal{P}_s|} \sum_{i \in \mathcal{P}_s} \ell_i.
\end{equation}
对整个 batch 的损失：
\begin{equation}
\mathcal{L}_{\mathrm{batch}} = \frac{1}{N} \sum_{s=1}^{N} \mathcal{L}_s.
\end{equation}

\subsubsection{训练过程：反向传播}

梯度只流过 $W_{\mathrm{proj}}$（$f_j$ 和 $q_s$ 均冻结）。链式法则路径：
\[
\mathcal{L} \xrightarrow{\partial/\partial z_j} z_j
\xrightarrow{\partial/\partial r_j} r_j^{(s)}
\xrightarrow{\partial/\partial v_j} v_j
\xrightarrow{\partial/\partial W_{\mathrm{proj}}} W_{\mathrm{proj}}.
\]

关键梯度：
\begin{equation}\label{eq:grad}
\frac{\partial \ell_i}{\partial z_j} = p_j - \mathbf{1}[j = i].
\end{equation}

\paragraph{推导过程：}
\begin{align}
\frac{\partial \ell_i}{\partial z_j}
&= \frac{\partial}{\partial z_j}\Big(-z_i + \log \sum_u e^{z_u}\Big) \notag\\
&= -\mathbf{1}[j=i] + \frac{e^{z_j}}{\sum_u e^{z_u}} = p_j - \mathbf{1}[j=i].
\end{align}

\paragraph{梯度的物理意义：}
\begin{itemize}[nosep]
    \item $j=i$（正样本）：梯度 $= p_i - 1 < 0$，SGD 会\textbf{增大} $z_i$，即增大 $q_s$ 与正样本 token 的余弦相似度（拉近）。
    \item $j \neq i$（负样本）：梯度 $= p_j > 0$，SGD 会\textbf{减小} $z_j$，即减小 $q_s$ 与负样本 token 的余弦相似度（推远）。
\end{itemize}
与 MoCo/SimCLR/CLIP 的实现形式一致（logits$/\tau$ 送入交叉熵）。

\subsubsection{互信息下界}

根据 CPC（Oord et al., 2018）：
\begin{equation}
I(x; c) \geq \log |\mathcal{C}_s| - \mathcal{L}_s.
\end{equation}
最优打分函数满足 $f^*(x,c) \propto p(x|c)/p(x)$（密度比）。训练 $W_{\mathrm{proj}}$ 即在冻结特征上学习该密度比的线性近似。

\begin{remark}[$|\mathcal{P}_s|=0$ 的处理]
当无正样本时损失未定义（除零）。策略：(1) 从批次中跳过该句；(2) 将最高 IoU 的 token 提升为正样本；(3) 仅用作其它句子的硬负样本。须记录选择并报告 $|\mathcal{P}_s|$ 分布。
\end{remark}

\begin{remark}[温度与空间先验交互]
空间先验在 logit 空间被 $1/\tau$ 放大。$\tau=0.07$，$\lambda_{\mathrm{spatial}}=0.3$ 时最大 logit 增益 $= 0.3/0.07 \approx 4.3$，非常显著。强烈建议做 $(\lambda_{\mathrm{spatial}}, \tau)$ 联合敏感性网格。
\end{remark}


%=================================================================
\section{M3：逐句 Token 门控生成}
%=================================================================

\subsection{Stage 3a：证据规划}

规划器接收粗 token 子集（$B_{\mathrm{plan}} = \min(32, B/4)$ 个 token），生成结构化计划：有序的句子主题列表，每句附带解剖关键词与目标深度层级。

\paragraph{输入 / 输出 / 参数来源。}
\begin{itemize}[nosep]
    \item \textbf{输入：}Token Bank 的前 $B_{\mathrm{plan}}$ 个 token（按 $S_i$ 降序选取），包含每个 token 的 bbox 中心坐标、深度层级和 volume 级元数据。LLM 仅接收文本化的 token 摘要（坐标、层级），不接收视觉特征向量。
    \item \textbf{输出：}结构化计划——一个有序列表 $\{(s_t, \mathrm{anatomy\_keyword}_t, [\ell_{\min}, \ell_{\max}]_t)\}_{t=1}^{T}$，其中 $s_t$ 为句子 topic，$\mathrm{anatomy\_keyword}_t$ 用于查询解剖 bbox，$[\ell_{\min}, \ell_{\max}]_t$ 为预期八叉树深度范围。$T$ 为规划句数（由 LLM 自动确定）。
    \item \textbf{参数来源：}$B_{\mathrm{plan}} = \min(32, B/4)$ — 限制规划 token 数以降低 LLM prompt 长度；LLM 为 Llama-3.1-8B-Instruct，temperature=0.3。
\end{itemize}

\subsection{Stage 3b：逐句路由}

对每个计划的句子主题 $s$：
\begin{equation}
E_s = \operatorname{top\text{-}k}_{i \in \{1,\ldots,B\}} \hat{r}_i^{(s)}, \quad |E_s| = k.
\end{equation}

\subsection{Stage 3c：Token 门控 LLM 生成}

LLM \textbf{逐句调用}。对第 $t$ 句：
\begin{itemize}
    \item \textbf{可见输入（视觉前缀）：}仅 $E_t$ 中 $k$ 个 token 嵌入（按空间顺序拼接）+ 此前已生成句子的文本历史。
    \item \textbf{输出：}单个报告句子 $y_t$。
\end{itemize}

\paragraph{输入 / 输出 / 参数来源（完整 prompt 结构）。}
\begin{itemize}[nosep]
    \item \textbf{输入组成：}(1) 系统消息（角色定义）；(2) Token block：对每个 $i \in E_t$ 提供 \texttt{token\_id}、\texttt{level}、\texttt{center (x,y,z)}、\texttt{bbox}、\texttt{volume}、\texttt{score}；(3) 证据卡（Evidence Card）：\texttt{laterality\_allowed}、\texttt{depth\_allowed}（见下文）；(4) 句子 topic $s_t$（来自 Stage 3a）；(5) 文本历史 $\{y_1, \ldots, y_{t-1}\}$（此前已生成的句子）。注意：LLM 接收的是 token 的\textbf{文本化空间描述}，而非原始视觉特征向量。
    \item \textbf{输出：}报告句子 $y_t$（自然语言文本），由 Llama-3.1-8B-Instruct 在 temperature=0.3 下生成。
    \item \textbf{参数来源：}temperature=0.3（较低以减少幻觉，同时保留一定多样性）；每句一次 LLM 调用。
\end{itemize}

\subsection{证据卡构建}

对 $E_s$ 中的每个 token，按 bbox 中心 $x$ 坐标分类侧别：
\begin{equation}
\mathrm{side}(i) = \begin{cases}
    \text{left}  & \text{if } x_c > x_{\mathrm{mid}} + \delta, \\
    \text{right} & \text{if } x_c < x_{\mathrm{mid}} - \delta, \\
    \text{cross} & \text{otherwise},
\end{cases}
\end{equation}
其中 $x_c = (x_{\min} + x_{\max})/2$，$x_{\mathrm{mid}} = W/2 = 64$（$128^3$ volume），$\delta = 0$（默认无容忍度）。

\textbf{主导侧比（Same-Side Ratio, SSR）：}
\begin{equation}
\rho = \frac{\max(n_L, n_R)}{n_L + n_R},
\end{equation}
其中 $n_L, n_R$ 为 left/right token 的数量（cross token 不计入分母）。

\paragraph{输入 / 输出 / 参数来源。}
\begin{itemize}[nosep]
    \item \textbf{输入：}$E_s$（$k$ 个 cited token 的 bbox 集合）+ volume 宽度 $W = 128$。
    \item \textbf{输出：}证据卡包含两个门控字段：
    \begin{itemize}[nosep]
        \item \texttt{laterality\_allowed}：\{left, right, bilateral, none\}。\textbf{v1}（B2'配置）：直接取主导侧。\textbf{v2}（B2'v2 配置，严格模式）：仅当 $\rho \geq 0.9$ 且非 cross token $\geq 2$ 时才允许 left/right 声明。
        \item \texttt{depth\_allowed}：\{in\_range, none\}。\textbf{v2}：仅当 $\geq 75\%$ 的 token 在预期深度范围 $[\ell_{\min}, \ell_{\max}]$ 内时允许深度相关描述。
    \end{itemize}
    \item \textbf{参数来源：}SSR 阈值 0.9 和最少 2 个非 cross token 的要求，来自 450-case 实验中 B2' → B2'v2 的消融发现——放松这些条件会导致 R1/R3 违规率上升。$\delta = 0$ 表示 token 中心恰好在中线上归为 cross。
\end{itemize}

LLM 在生成时遵循证据卡约束：若 \texttt{laterality\_allowed=none}，则生成文本中不包含侧别词（left/right）；若 \texttt{depth\_allowed=none}，则不包含深度定位词（upper/lower）。

\subsection{接口正确性：构造性引用}

\begin{definition}[构造性引用]
令 $E_t = \{e_{t,1},\ldots,e_{t,k}\}$ 为路由到第 $t$ 句的证据集，$\hat{E}_t$ 为输出中记录的引用集。由流水线构造：
\begin{equation}\label{eq:cite}
\hat{E}_t \equiv E_t \quad \text{（精确集合相等）}。
\end{equation}
\end{definition}

这保证了\textbf{溯源正确性}：每个引用都是模型可见的证据，每个可见 token 都被引用。无漏报，无多报。

\begin{remark}[接口正确性 $\neq$ 语义支持]
方程~\eqref{eq:cite} 保证引用溯源正确，但\emph{不}保证生成文本被引用证据语义支持。模型可能忽略视觉 token 而依赖语言先验（``虽引用正确但实质幻觉''）。语义支持须由 M4 校验器 + 外部 grounded 评估验证。
\end{remark}

\begin{remark}[跨句一致性]
多次独立 LLM 调用可能导致篇章漂移：同一病灶在不同句中描述不一致、侧别矛盾或冗余。文本历史前缀部分缓解此问题，但需显式跟踪跨句一致性指标。
\end{remark}


%=================================================================
\section{M4：规则校验器与一次性重路由（Stage 4--5）}
%=================================================================

\subsection{Stage 4：规则化校验}

对每句及其引用 token 的边界框应用五类规则：

\paragraph{R1 --- 侧别一致性。}
若句子提及``左''（``右''），则所有引用 token 的 bbox 中心须满足 $x_{\mathrm{center}} > x_{\mathrm{mid}}$（$< x_{\mathrm{mid}}$），其中 $x_{\mathrm{mid}}$ 为标准化坐标系（LPS/RAS）下的中线。

% >>> 新增：R1 工程增强
\paragraph{R1 实现增强。}
实际实现中，R1 引入三项改进以降低误报：
\begin{enumerate}[nosep]
    \item \textbf{Ratio 模式：}将 cited tokens 按 $\mathrm{token\_side}()$ 分为 same/opposite/cross 三类，仅当同侧比例 $\mathrm{same\_ratio} = |\mathrm{same}| / (|\mathrm{same}| + |\mathrm{opp}|) < \theta_{\mathrm{ratio}}$ 时触发 R1（默认 $\theta_{\mathrm{ratio}} = 0.6$），取代原有的 all-or-nothing 判定。中线附近 token（cross）排除在分母之外。
    \item \textbf{否定句豁免：}否定句（如``未见左侧积液''）cite 的是背景区域 token，不要求侧别一致，跳过 R1。
    \item \textbf{中线解剖词豁免：}mediastinum / trachea / carina / esophagus / aorta / spine / vertebra / sternum 等天然跨中线结构，cited tokens 必然覆盖双侧，跳过 R1。
\end{enumerate}

\paragraph{R2 --- 解剖区域一致性。}
引用 token 的 bbox 须与句子暗示的解剖区域有非平凡 IoU（$> \tau_{\mathrm{anatomy}}$）。

% >>> 新增：R2 工程增强
\paragraph{R2 实现增强。}
\begin{enumerate}[nosep]
    \item \textbf{支持率模式：}计算 cited tokens 中 $\IoU > \tau_{\mathrm{IoU}}$ 的比例 $\mathrm{good\_ratio}$，仅当 $\mathrm{good\_ratio} < \theta_{\mathrm{support}}$（默认 0.8）时触发 R2，避免因单个离群 token 导致整句违规。
    \item \textbf{大结构豁免：}bilateral / left lung / right lung 等覆盖大范围的解剖结构，其 bbox 与小 token bbox 的 IoU 天然极低，属于结构性误报，在 R2 中予以豁免（加入 \texttt{r2\_skip\_keywords}）。对于 mediastinum，通过将 $\tau_{\mathrm{IoU}}$ 从 0.05 降至 0.04（经 mediastinum 阈值 sweep 实验确定的临界点）来消除误报，而非直接豁免。
\end{enumerate}

\paragraph{R3 --- 深度层级一致性。}
引用 token 的八叉树深度须在所述发现类型的预期范围内。

\paragraph{R4 --- 大小/范围合理性。}
引用 token 的联合 bbox 须与所述发现的典型大小范围一致。

\paragraph{R5 --- 否定处理。}
若检测到否定（``未见\ldots''），校验逻辑取反：引用 token 应\emph{不}表现被否定的发现。

每条规则 $r$ 输出：二值判定（通过/失败）及严重度 $\mathrm{sev}_r \in [0,1]$。默认严重度配置：$\mathrm{sev}_{\mathrm{R1}}=1.0$、$\mathrm{sev}_{\mathrm{R2}}=0.8$、$\mathrm{sev}_{\mathrm{R3}}=0.7$、$\mathrm{sev}_{\mathrm{R4}}=0.6$、$\mathrm{sev}_{\mathrm{R5}}=1.0$。

\paragraph{M4 整体输入 / 输出。}
\begin{itemize}[nosep]
    \item \textbf{输入（每句）：}(1) 生成的句子 $y_t$（文本）；(2) cited token 集合 $E_t$（每个 token 的 bbox、depth level）；(3) 解剖 bbox $\mathrm{bbox}_{\mathrm{anatomy}}^{(s)}$（来自 \texttt{DEFAULT\_ANATOMY\_BOXES} 规则表）；(4) volume 宽度 $W=128$（用于计算 $x_{\mathrm{mid}}$）；(5) 同一 volume 的所有已生成句子（R6 跨句检查需要）。
    \item \textbf{输出（每句）：}违规列表 $\{(r, \mathrm{sev}_r, \mathrm{details})\}$，其中 $r \in \{\mathrm{R1}, \mathrm{R2}, \mathrm{R3}, \mathrm{R6a}, \mathrm{R6b}\}$。无违规则为空列表。
    \item \textbf{参数来源：}$\theta_{\mathrm{ratio}} = 0.6$（R1）和 $\theta_{\mathrm{support}} = 0.8$（R2）由 450-case 实验的阈值 sweep 确定。$\tau_{\mathrm{IoU}} = 0.04$（mediastinum）/ $0.05$（其它）由解剖结构大小决定。严重度权重为预设值：R1（侧别）和 R5（否定）为 1.0（最严重），R3（深度）为 0.7（允许一定宽容度）。
\end{itemize}

\begin{remark}[当前实验中的规则启用状态]
450/450 实验中：R4（大小/范围合理性）\textbf{已禁用}（阈值未校准）；R5 的回退词典\textbf{已禁用}（仅保留元数据级否定检测）。R1--R3 和 R6 正常启用。
\end{remark}

% >>> 新增：R6 跨句一致性
\paragraph{R6 --- 跨句一致性（增强规则）。}
对同一解剖关键词下的所有句子进行跨句检查：
\begin{itemize}[nosep]
    \item \textbf{R6a --- 侧别冲突：}若同一解剖词在不同句子中声明相反侧别（一句说``左''，另一句说``右''），标记冲突（$\mathrm{sev}=0.7$）。
    \item \textbf{R6b --- 存在/缺失冲突：}若同一解剖词在一句中为阳性发现、在另一句中为否定发现，标记冲突（$\mathrm{sev}=0.8$）。
\end{itemize}

\subsection{Stage 5：句级软惩罚与重路由}

由于全局综合重要性分数 $S_i$ 不直接决定单句的 Token 选取，惩罚应该尝试去直接施加于跨模态路由分数 $\hat{r}_i^{(s)}$ 以改变 top-$k$ 采样结果。

\begin{definition}[句级严重度]
对第 $t$ 句（对应主题 $s$）涉及违规的 token $i$，汇聚其触发的规则严重度：
\begin{equation}
\mathrm{sev}_i^{(s)} = \max_{r \in \text{触发的规则}} \mathrm{sev}_r(i, s) \in [0,1].
\end{equation}
（注：若 token $i$ 未违规，则 $\mathrm{sev}_i^{(s)} = 0$）。
\end{definition}

\begin{definition}[对数平滑惩罚]
为避免硬性阈值截断导致的上下文断裂，对 token $i$ 的路由分数施加基于严重度的对数平滑惩罚：
\begin{equation}\label{eq:penalty}
\hat{r}_i^{(s)\prime} = \hat{r}_i^{(s)} - \gamma \cdot \ln(1 + \mathrm{sev}_i^{(s)}),
\end{equation}
其中 $\gamma > 0$ 为惩罚缩放因子（默认取 $\gamma = 2.0$）。

\paragraph{输入 / 输出 / 参数来源。}
\begin{itemize}[nosep]
    \item \textbf{输入：}$\hat{r}_i^{(s)}$（原始路由分数，来自公式~\ref{eq:spatial} 或 \ref{eq:anat_primary}）+ $\mathrm{sev}_i^{(s)}$（该 token 在该句的违规严重度汇聚值）。
    \item \textbf{输出：}$\hat{r}_i^{(s)\prime}$ — 惩罚后的路由分数。对所有 $B$ 个 token 重新排序取 top-$k$，得到新证据集 $E_s'$，用于重新生成句子。
    \item \textbf{参数来源：}$\gamma = 2.0$ 确保最大惩罚 $\gamma \ln 2 \approx 1.386$ 足以将严重违规 token（$\mathrm{sev}=1$）从 top-$k$ 中挤出——因为正常路由分数范围 $[-1, 1.3]$（含 $\lambda_{\mathrm{spatial}}$），惩罚量占满量程的约 $50\%$。$\ln(1+x)$ 的对数形式使轻微违规（$\mathrm{sev} \ll 1$）受到较小惩罚，避免过度修正。
\end{itemize}
\end{definition}

\begin{proposition}[惩罚有界性与保序性]
(i) 已知 $\hat{r}_i^{(s)} \in [-1, 1+\lambda_{\text{spatial}}]$ 且 $\mathrm{sev}_i^{(s)} \in [0,1]$。新路由分数严格有界：
\begin{equation}
\hat{r}_i^{(s)\prime} \in [-1 - \gamma \ln 2, 1+\lambda_{\text{spatial}}].
\end{equation}
若取默认值 $\gamma=2.0$，最大惩罚量 $\Delta \approx 1.386$，足以抵消最大可能的空间先验增益（$\lambda_{\text{spatial}}$）并将其挤出 top-$k$。\\
(ii) 对于未违规的合规 token 集合（$\mathrm{sev}_i^{(s)} = 0$），其相对大小排序保持严格不变。
\end{proposition}

% >>> 新增：LLM 语义裁判
\begin{remark}[LLM 语义裁判（Stage 5 增强）]\label{rmk:llm_judge}
在对数平滑惩罚之前，可选地引入 LLM 语义裁判对 Stage 4 的每条违规做二次确认。LLM 接收违规句文本和规则描述，输出 $\{\mathrm{confirmed}: \mathrm{bool},\, \mathrm{severity}: [0,1],\, \mathrm{reasoning}: \mathrm{str}\}$。仅对 LLM 确认的违规执行惩罚和重路由，未确认的违规视为误报并移除。该机制弥补了纯空间/几何规则缺乏语义理解能力的局限。当前实验使用 Llama-3.1-8B-Instruct（本地 bfloat16 推理），在 450/450 实验中过滤了约 54\% 的 Stage 4 违规（从 974 条降至 449 条确认违规）。支持 Ollama / OpenAI / Anthropic / HuggingFace 四种后端。
\end{remark}

\paragraph{重路由协议：}
\begin{enumerate}
    \item 用更新后的路由分数 $\{\hat{r}_i^{(s)\prime}\}$ 重新执行 top-$k$ 采样，获取新证据集 $E_s'$。
    \item 用新证据集 $E_s'$ 重新生成失败的句子。
    \item 对重新生成的句子再次执行 Stage 4 校验。
    \item 若仍失败：\emph{降格表述}（\texttt{despecify\_text()}：移除空间定位词，包括侧别词 left/right、方位词 upper/lower/anterior/posterior、层级词 apical/basal/hilar 等），然后用降格后的主题重新生成。
\end{enumerate}
仅一次重试，无递归重路由。


%=================================================================
\section{M5：匹配计算量评估与统计协议}
%=================================================================


\subsection{计算量匹配}

\begin{definition}[LLM Token 代价]
对 $J$ 次 LLM 调用：
\begin{equation}\label{eq:cllm}
C_{\mathrm{LLM}} = \sum_{j=1}^{J} (|p_j| + |g_j|),
\end{equation}
其中 $|p_j|$ = 提示长度，$|g_j|$ = 生成长度。

\paragraph{输入 / 输出 / 参数来源。}
\begin{itemize}[nosep]
    \item \textbf{输入：}所有 LLM 调用的 prompt token 数 $|p_j|$ 和生成 token 数 $|g_j|$。包括 Stage 3a（规划）、Stage 3c（生成）、Stage 5a（LLM 裁判）的每次调用。
    \item \textbf{输出：}$C_{\mathrm{LLM}}$（总 token 代价）。用于配置间的公平计算量对比——例如 C2'（含裁判）比 B2'v2（无裁判）多消耗 $\sim 3\%$ 的 LLM token。
    \item \textbf{当前实现简化：}代码中按调用次数统计 $C_{\mathrm{LLM}} = n_{\mathrm{gen}} + n_{\mathrm{judge}}$，尚未逐调用记录 $|p_j|, |g_j|$。Table 2 中的 $C_{\mathrm{LLM}}$ 列即为调用次数。
\end{itemize}
\end{definition}

\begin{remark}[$C_{\mathrm{LLM}}$ 实现简化]
当前代码按\textbf{调用次数}统计 $C_{\mathrm{LLM}} = n_{\mathrm{gen}} + n_{\mathrm{judge}}$（包括重路由和降格表述产生的额外调用），尚未按公式逐调用记录 prompt/生成 token 数。完整的逐 token 统计需在 \texttt{stage3c\_generator} 中记录 $|p_j|$ 和 $|g_j|$。
\end{remark}

\subsection{注意力代价代理}

\subsubsection{朴素（无 KV-Cache）代理}
\begin{equation}\label{eq:naive}
C_{\mathrm{attn}}^{\mathrm{naive}} = \sum_{j=1}^{J} (|p_j| + |g_j|)^2.
\end{equation}

\subsubsection{KV-Cache 代理：逐步推导}

对第 $j$ 次调用，提示长度 $p = |p_j|$，生成长度 $g = |g_j|$：

\paragraph{第一步：预填充。} 对提示做完整自注意力：代价 $\propto p^2$。

\paragraph{第二步：增量生成。} 在生成第 $t$ 步（$t=1,\ldots,g$），新 token 对长度为 $p+(t-1)$ 的上下文做注意力：
\begin{equation}
\text{第 } t \text{ 步代价} = p + (t-1).
\end{equation}

\paragraph{第三步：对生成步求和。}
\begin{align}
\sum_{t=1}^{g}(p + (t-1))
&= gp + \sum_{t=0}^{g-1} t = gp + \frac{g(g-1)}{2}. \label{eq:gensum}
\end{align}

\paragraph{第四步：总代理。}
\begin{equation}\label{eq:cache}
\boxed{
C_{\mathrm{attn}}^{\mathrm{cache}} = \sum_{j=1}^{J}\left(
    |p_j|^2 + |g_j|\cdot|p_j| + \frac{|g_j|(|g_j|-1)}{2}
\right).
}
\end{equation}

\begin{remark}[替代约定]
若每步 $t$ 对长度 $p+t$ 的上下文做注意力（含当前 token）：
\[
\sum_{t=1}^{g}(p+t) = gp + \frac{g(g+1)}{2}.
\]
差异为 $O(g)$，作为代理可接受，但须明确约定。
\end{remark}


\subsection{病人级配对 Bootstrap 置信区间}

\begin{definition}[配对 Bootstrap CI]
给定 $n$ 个病人的指标对 $\{(m_i^A, m_i^B)\}_{i=1}^{n}$，令 $\delta_i = m_i^A - m_i^B$：
\begin{enumerate}
    \item 对 $b = 1,\ldots,R$（默认 $R=5000$）：
    \begin{enumerate}[label=(\alph*)]
        \item 有放回抽 $n$ 个下标。
        \item 计算 $\bar{\delta}^{*(b)} = \frac{1}{n}\sum_{l=1}^n \delta_{i_l^*}$。
    \end{enumerate}
    \item 对 $R$ 个 bootstrap 均值排序。
    \item $(1-\alpha)$ 分位数 CI：
    \begin{equation}
    \mathrm{CI}_{1-\alpha} = \Big[\bar{\delta}^*_{(\lfloor R\alpha/2\rfloor)},\;\bar{\delta}^*_{(\lceil R(1-\alpha/2)\rceil)}\Big].
    \end{equation}
\end{enumerate}
\end{definition}

\paragraph{输入 / 输出 / 参数来源。}
\begin{itemize}[nosep]
    \item \textbf{输入：}$n$ 个病人的指标对 $\{(m_i^A, m_i^B)\}$，其中 $A, B$ 为两个待比较的消融配置。$m_i$ 可以是违规率、BLEU-4、CF 等任意标量指标，按病人聚合（每个病人有多个句子，先取均值）。
    \item \textbf{输出：}$95\%$ 置信区间 $[\bar{\delta}^*_{\mathrm{lo}}, \bar{\delta}^*_{\mathrm{hi}}]$。若 CI 不跨零则配置间差异显著。
    \item \textbf{参数来源：}$R=5000$（bootstrap 重采样次数）为统计学标准推荐值（1000--10000 范围内），在精度和计算量之间取平衡。$\alpha=0.05$（$95\%$ CI）为惯例。
\end{itemize}

\paragraph{为何按病人级采样？}同一 CT 扫描内的句子相关（相同解剖、病理、图像质量）。按句级采样会低估方差，产生虚假窄 CI。病人级重采样保留组内相关性。


\subsection{Holm 步降多重比较校正}

\begin{definition}[Holm 校正]
给定 $m$ 个假设检验，按 $p$ 值排序 $p_{(1)} \leq \cdots \leq p_{(m)}$：
\begin{enumerate}
    \item 对 $k = 1,2,\ldots,m$：比较 $p_{(k)}$ 与 $\alpha/(m-k+1)$。
    \item 若 $p_{(k)} \leq \alpha/(m-k+1)$ 且此前全部被拒绝，则拒绝 $H_{(k)}$。
    \item 在第一个未被拒绝处停止。
\end{enumerate}
控制族错误率（FWER）于 $\alpha$ 水平；一致优于 Bonferroni。
\end{definition}

同时报告校正后与未校正 $p$ 值，确保透明度。

\begin{remark}[M5 实现状态]
Bootstrap CI（$R=5000$，按病人级抽样）和 Holm 步降校正均已在 \texttt{analyze\_outputs.py} 中实现（\texttt{\_bootstrap\_ci()} 和 \texttt{\_holm\_correction()}），可在 450/450 实验结果上直接运行。
\end{remark}


%=================================================================
\section{超参数总结}
%=================================================================

% >>> 修改：更新推荐值、新增 γ/ε/θ 行
\begin{longtable}{p{3.6cm}cp{3cm}cp{3.6cm}}
\caption{默认值、文献范围与消融网格。}\label{tab:hyper}\\
\toprule
\textbf{超参数} & \textbf{默认} & \textbf{文献范围} & \textbf{推荐} & \textbf{消融网格} \\
\midrule\endfirsthead
\toprule
\textbf{超参数} & \textbf{默认} & \textbf{文献范围} & \textbf{推荐} & \textbf{消融网格} \\
\midrule\endhead

Token bank 大小 $B$ & 64 & Flamingo $\leq 64$; BLIP-2 $=32$ & 128 & $\{32,64,128\}$ \\
每句 $k$ & 8 & ContextCite top-$k$ & 8 & $\{4,8,16\}$ \\
规划 token 数 $B_{\mathrm{plan}}$ & $\min(32,B/4)$ & BLIP-2 32 queries & 32 & $\{0,16,32\}$ \\
伪影阈值 $\tau_A$ & 0.85 & 临床 SNR/CNR & 0.85 & $\{0.75,0.85,0.90\}$ \\
门控陡峭度 $k_{\mathrm{gate}}$ & 15 & 工程选择 & 15 & $\{10,15,20\}$ \\
权重 $(w_e,w_v,w_b)$ & $(0.5,0.3,0.2)$ & 多信号融合 & 默认 & $\pm 0.1$ 扰动 \\
评分融合 $\lambda_h$ & 0.6 & 无统一标准 & 0.6 & $\{0.4,0.6,0.8\}$ \\
边界混合 $\beta$ & 0.1 & 工程选择 & 0.1 & $\{0,0.05,0.1,0.2\}$ \\
空间先验 $\lambda_{\mathrm{spatial}}$ & 0.3 & 实验消融网格 & 0.3 & $\{0,0.1,0.3,0.5\}$ \\
IoU 阈值 $\tau_{\mathrm{IoU}}$ & 0.1 & 弱监督标准 & 0.04 & $\{0.03,0.04,0.05,0.1\}$ \\
InfoNCE 温度 $\tau$ & 0.07 & MoCo 0.07; SimCLR 0.1 & 0.07 & $\{0.05,0.07,0.1,0.2\}$ \\
惩罚缩放 $\gamma$ & 2.0 & 工程选择 & 2.0 & $\{1.0,2.0,3.0\}$ \\
Anatomy tiebreak $\epsilon$ & 0.05 & 工程选择 & 0.05 & $\{0.01,0.05,0.1\}$ \\
R1 ratio 阈值 $\theta_{\mathrm{ratio}}$ & 1.0 & 工程选择 & 0.6 & $\{0.4,0.6,0.8,1.0\}$ \\
R2 支持率 $\theta_{\mathrm{support}}$ & 1.0 & 工程选择 & 0.8 & $\{0.6,0.8,1.0\}$ \\
Holm 校正 & 是 & 标准 FWER & 是 & 固定 \\
Bootstrap $R$ & 5000 & 1k--10k & 5000 & $\{1000,5000,10000\}$ \\
\bottomrule
\end{longtable}


%=================================================================
\section{未决问题清单}
%=================================================================

% >>> 修改：已在代码中解决的标记 \checkmark，未解决的保留 □
\begin{enumerate}[label=$\square$]
    \item[\checkmark] \textbf{$\normop(\cdot)$ 与 $q(\cdot)$ 的形式化定义：}代码采用 min-max 归一化（$\normop$）和分位数排名（$q$），按病例/深度层独立计算。Ties 按原始 enumerate 顺序分配排名。
    \item[\checkmark] \textbf{3D NMS 规格：}轴对齐 3D IoU，可配置 NMS 阈值（默认 1.0 即不做 NMS），全局执行。
    \item[\checkmark] \textbf{侧别坐标系：}预处理中统一为 LPS 朝向（\texttt{sitk.DICOMOrient(img, "LPS")}），x 轴：$0=$ 患者右侧，$x_{\max}=$ 患者左侧。
    \item[$\square$] \textbf{$|\mathcal{P}_s|=0$ 处理：}当前实现抛出异常；尚未报告 $|\mathcal{P}_s|$ 分布。
    \item[$\square$] \textbf{$\lambda_{\mathrm{spatial}} \times \tau$ 联合敏感性：}提供网格消融，防止空间项在小 $\tau$ 下压倒语义路由。
    \item[\checkmark] \textbf{$\mathrm{sev}_i$ 归一化：}各规则严重度在配置中预设 $\in [0,1]$；R1--R5 到逐 token 的映射通过 \texttt{severity\_by\_rule} 字典实现。
    \item[\checkmark] \textbf{跨句一致性：}R6a（侧别冲突）和 R6b（存在/缺失冲突）已实现。
    \item[$\square$] \textbf{注意力代理约定：}当前代码使用 $k \cdot B$ 简化代理，尚未按公式~\eqref{eq:cache} 逐调用计算。
\end{enumerate}


%=================================================================
\section{方法--相关工作交叉引用}
%=================================================================

\begin{longtable}{p{1cm}p{3cm}p{7cm}p{2.5cm}}
\caption{方法 $\times$ 相关工作对照。}\label{tab:related}\\
\toprule
\textbf{方法} & \textbf{工作} & \textbf{对应关系} & \textbf{出处位置} \\
\midrule\endfirsthead
\toprule
\textbf{方法} & \textbf{工作} & \textbf{对应关系} & \textbf{出处位置} \\
\midrule\endhead

M1 & OctNet & 稀疏分层空间划分；叶节点 pooled 特征 & 摘要+引言 \\
M1 & TokenLearner & 少量自适应 token 降低 $O(n^2)$ 注意力代价 & Sec 3 \\
M2 & CLIP & 余弦相似+温度缩放 softmax & Sec 3.1 \\
M2 & MoCo & $\tau=0.07$ 默认；logits/$\tau$+CE 实现 & Sec 3.3 \\
M2 & CPC（InfoNCE） & 分类交叉熵解释；密度比；MI 下界 & Eq (4)--(5) \\
M2 & CT-GLIP & 3D grounded VLP；$\tau=0.07$；局部对齐 & 方法部分 \\
M3 & RECLAIM & 逐句引用；约束解码 & Sec 4 \\
M3 & LongCite/CoF & 逐句 top-$l$ 检索+引用抽取 & Sec 4.1 \\
M3/4 & ContextCite & top-$k$ 归因优于全上下文验证 & Fig 5 \\
M3/4 & CoVe & 规划$\to$核查$\to$修订结构 & 摘要+Sec 3 \\
M4 & 侧别错误文献 & 关键词/正则的侧别筛查 & 综述 \\
M5 & Bootstrap 综述 & 重采样分位数 CI & 方法部分 \\
M5 & Holm 校正 & 步降 FWER 控制 & 公式描述 \\
\bottomrule
\end{longtable}

%=================================================================
% 实验结果表格
%=================================================================
\newpage
\section{实验结果}

以下表格基于 5K 数据集的 test split（CT-RATE 250 例 + RadGenome-ChestCT 250 例，共 496 个病例、3{,}979 个句子）。外部基线数据引用自各方法的原始论文。

\paragraph{消融链配置说明。}
消融链按以下顺序逐步叠加模块，每一步仅增加一个组件：
\begin{itemize}[nosep]
    \item \textbf{A0}：基线。$W_{\mathrm{proj}}=I$（单位矩阵），纯空间 IoU 路由，不生成句子（直接使用原始 topic）。
    \item \textbf{A1}：在 A0 基础上训练 $W_{\mathrm{proj}}$（InfoNCE），启用语义 + 空间融合路由。
    \item \textbf{E1}：在 A1 基础上加入空间过滤（spatial filter）+ 语义重排（semantic rerank）。
    \item \textbf{B2$'$}：在 E1 基础上加入 Evidence Card v1（SSR 侧别门控 + 深度门控）+ LLM 生成句子。
    \item \textbf{B2$'$v2}：将 Evidence Card 升级为 v2（更严格的门控条件，不满足时截断 token）。
    \item \textbf{C2$'$}：在 B2$'$ 基础上加入 LLM Judge（Stage 5），对高违规句子进行 LLM 二次审核。
    \item \textbf{D2}：在 C2$'$ 基础上加入修复执行器（reroute / despecify）。
\end{itemize}

%─────────────────────────────────────────────────────────
\subsection{Table 1：与外部基线的 NLG 质量对比（主表）}
%─────────────────────────────────────────────────────────

\begin{table}[H]
\centering
\caption{与外部基线的报告生成质量对比。基线结果引用自各方法原始论文；ProveTok 使用 5K 划分（\texttt{seed=42}，2000/250/250）。}\label{tab:main_comparison}
\footnotesize
\setlength{\tabcolsep}{4pt}
\begin{tabular}{@{}l l l r r r@{}}
\toprule
\textbf{Dataset} & \textbf{Method} & \textbf{Protocol} & \textbf{B-4}$\uparrow$ & \textbf{MTR}$\uparrow$ & \textbf{R-L}$\uparrow$ \\
\midrule
CT-RATE   & CT2Rep      & published & .172 & .173 & .243 \\
CT-RATE   & CT-AGRG     & published & .172 & .196 & .280 \\
CT-RATE   & \textbf{ProveTok} & ours & \textbf{.467} & \textbf{.603} & \textbf{.626} \\
\midrule
RadGenome & MedVInT     & published & .246 & .404 & .326 \\
RadGenome & Reg2RG      & published & .249 & .441 & .367 \\
RadGenome & \textbf{ProveTok} & ours & \textbf{.506} & \textbf{.660} & \textbf{.680} \\
\bottomrule
\end{tabular}
\\[4pt]
{\scriptsize B-4 = BLEU-4; MTR = METEOR; R-L = ROUGE-L.}
\end{table}

%─────────────────────────────────────────────────────────
\subsection{Table 2：消融链——空间一致性与计算开销}
%─────────────────────────────────────────────────────────

\begin{table}[H]
\centering
\caption{消融链（5K test split，496 例，3{,}979 句）。Viol.\% = 句子级违规率；R1/R3/R6b = 各规则违规计数；$C_{\mathrm{LLM}}$ = LLM 调用次数；$N_{\mathrm{reroute}}$ = 重路由次数。}\label{tab:ablation_main}
\small
\begin{tabular}{@{}l l r r r r r r r@{}}
\toprule
\textbf{Config} & \textbf{Description} & \textbf{Viol.\%}$\downarrow$ & \textbf{R1} & \textbf{R3} & \textbf{R6b} & \textbf{Total} & $C_{\mathrm{LLM}}$ & $N_{\mathrm{reroute}}$ \\
\midrule
A0     & Identity $W$ + Spatial            & 7.59 &   0 & 237 & 65 & 302 &    0 &   0 \\
A1     & Trained $W$ + Spatial             & 6.01 &   0 & 174 & 65 & 239 &    0 &   0 \\
E1     & Spatial filter + Semantic rerank  & 6.18 &  95 &  86 & 65 & 246 &    0 &   0 \\
B2$'$  & + Evidence Card v1               & 6.01 & 109 &  86 & 44 & 239 & 3979 &   0 \\
\textbf{B2$'$v2} & \textbf{+ Evidence Card v2} & \textbf{4.25} & \textbf{98} & \textbf{23} & \textbf{48} & \textbf{169} & \textbf{3979} & \textbf{0} \\
C2$'$  & + LLM Judge (Stage 5)            & 6.18 & 113 &  86 & 47 & 246 & 4112 & 128 \\
D2     & + Repair executor                & 5.96 & 102 &  86 & 49 & 237 & 4116 & 131 \\
\bottomrule
\end{tabular}
\end{table}

%─────────────────────────────────────────────────────────
\subsection{Table 3：Grounding 质量指标}
%─────────────────────────────────────────────────────────

\begin{table}[H]
\centering
\caption{Grounding 指标（5K test split）。Overlap = 均值重叠率；Hit@$k$ = top-$k$ 命中率；RP = Routing Precision；Lat.Acc = 侧别准确率；CF = Citation Faithfulness；DF = Depth Faithfulness。``---'' = 无 LLM 生成，无法计算。}\label{tab:grounding}
\small
\begin{tabular}{@{}l r r r r r r r r@{}}
\toprule
\textbf{Config} & \textbf{Overlap} & \textbf{Hit@1} & \textbf{Hit@4} & \textbf{Hit@8} & \textbf{RP} & \textbf{Lat.Acc} & \textbf{CF} & \textbf{DF} \\
\midrule
A0     & 0.931 & 100.0 & 100.0 & 100.0 & 100.0 & 100.0 & ---  & ---  \\
A1     & 0.931 & 100.0 & 100.0 & 100.0 & 100.0 & 100.0 & ---  & ---  \\
E1     & 0.685 &  76.8 &  96.1 &  99.3 &  75.4 &  97.2 & ---  & ---  \\
B2$'$  & 0.685 &  76.8 &  96.1 &  99.3 &  75.4 &  96.9 & 76.9 & 97.3 \\
\textbf{B2$'$v2} & \textbf{0.689} & \textbf{76.8} & \textbf{95.9} & \textbf{98.2} & \textbf{75.8} & \textbf{97.1} & \textbf{77.6} & \textbf{98.4} \\
C2$'$  & 0.685 &  76.8 &  96.1 &  99.3 &  75.4 &  96.7 & 81.1 & 97.3 \\
D2     & 0.685 &  76.8 &  96.1 &  99.3 &  75.4 &  97.0 & 80.6 & 97.3 \\
\bottomrule
\end{tabular}
\end{table}

%─────────────────────────────────────────────────────────
\subsection{Table 4：消融链 NLG 质量}
%─────────────────────────────────────────────────────────

\begin{table}[H]
\centering
\caption{消融链 NLG 质量（5K test split）。仅有 LLM 生成的配置（B2$'$ 及之后）可计算。}\label{tab:nlg}
\small
\begin{tabular}{@{}l r r r@{}}
\toprule
\textbf{Config} & \textbf{BLEU-4} & \textbf{ROUGE-L} & \textbf{METEOR} \\
\midrule
B2$'$            & 0.4707 & 0.6425 & 0.6175 \\
\textbf{B2$'$v2} & \textbf{0.4725} & \textbf{0.6426} & \textbf{0.6190} \\
C2$'$            & 0.4771 & 0.6431 & 0.6218 \\
D2               & 0.4738 & 0.6434 & 0.6217 \\
\bottomrule
\end{tabular}
\end{table}

%─────────────────────────────────────────────────────────
\subsection{Table 5：统计显著性检验}
%─────────────────────────────────────────────────────────

\begin{table}[H]
\centering
\caption{Bootstrap 95\% CI（$R=5{,}000$，病人级重采样）。}\label{tab:bootstrap_ci}
\small
\begin{tabular}{@{}l l r r r r@{}}
\toprule
\textbf{Config} & \textbf{Description} & \textbf{Viol.\%} & \textbf{CI$_{\mathrm{lo}}$} & \textbf{CI$_{\mathrm{hi}}$} & $n$ \\
\midrule
A0     & Identity $W$ + Spatial   & 7.59 & 6.84 & 8.35 & 496 \\
A1     & Trained $W$ + Spatial    & 5.96 & 5.27 & 6.64 & 496 \\
E1     & + Spatial filter + Rerank & 6.04 & 5.34 & 6.77 & 496 \\
B2$'$  & + Evidence Card v1       & 5.85 & 5.14 & 6.60 & 496 \\
\textbf{B2$'$v2} & \textbf{+ Evidence Card v2} & \textbf{4.19} & \textbf{3.60} & \textbf{4.83} & \textbf{496} \\
C2$'$  & + LLM Judge              & 6.00 & 5.25 & 6.79 & 496 \\
D2     & + Repair executor        & 5.79 & 5.07 & 6.55 & 496 \\
\bottomrule
\end{tabular}
\end{table}

\begin{table}[H]
\centering
\caption{相邻配置 Permutation Test + Holm 校正（$R=5{,}000$，$\alpha=0.05$）。$\Delta$ = Viol.\% 差值（百分点）；$^*$ = Holm 校正后显著。}\label{tab:pairwise}
\small
\begin{tabular}{@{}l r r r r r l@{}}
\toprule
\textbf{Comparison} & \textbf{Mean$_A$} & \textbf{Mean$_B$} & $\Delta$ & $p$ & $p_{\mathrm{adj}}$ & \textbf{Sig.} \\
\midrule
A0 $\to$ A1         & 7.59 & 5.96 & $-$1.63 & 0.0020 & 0.0080 & $^*$ \\
A1 $\to$ E1         & 5.96 & 6.04 & $+$0.08 & 0.8702 & 1.0000 &  \\
E1 $\to$ B2$'$      & 6.04 & 5.85 & $-$0.19 & 0.7333 & 1.0000 &  \\
B2$'$ $\to$ B2$'$v2 & 5.85 & 4.19 & $-$1.65 & 0.0010 & 0.0060 & $^*$ \\
B2$'$v2 $\to$ C2$'$ & 4.19 & 6.00 & $+$1.80 & 0.0010 & 0.0060 & $^*$ \\
C2$'$ $\to$ D2      & 6.00 & 5.79 & $-$0.21 & 0.7059 & 1.0000 &  \\
\bottomrule
\end{tabular}
\end{table}

%─────────────────────────────────────────────────────────
\subsection{Table 6：错误分析——按解剖区域的违规分布（B2$'$v2）}
%─────────────────────────────────────────────────────────

\begin{table}[H]
\centering
\caption{B2$'$v2 配置下按解剖关键词的违规分布。仅列出有违规的解剖区域。}\label{tab:anatomy_viol}
\small
\begin{tabular}{@{}l r r r r r r@{}}
\toprule
\textbf{Anatomy} & \textbf{Sentences} & \textbf{Violated} & \textbf{Viol.\%} & \textbf{R1} & \textbf{R3} & \textbf{R6b} \\
\midrule
bilateral    & 1122 & 67 &  6.0 &  0 & 23 & 45 \\
right lung   &  198 & 58 & 29.3 & 56 &  0 &  3 \\
left lung    &  131 & 36 & 27.5 & 36 &  0 &  0 \\
(no anatomy) & 2293 &  6 &  0.3 &  6 &  0 &  0 \\
\bottomrule
\end{tabular}
\end{table}

%─────────────────────────────────────────────────────────
\subsection{主要发现}
%─────────────────────────────────────────────────────────

\begin{enumerate}[nosep]
    \item \textbf{ProveTok 大幅超越外部基线}：在 CT-RATE 上 BLEU-4 = 0.467 vs CT2Rep/CT-AGRG 的 0.172（$+$171\%）；在 RadGenome-ChestCT 上 BLEU-4 = 0.506 vs Reg2RG 的 0.249（$+$103\%）。METEOR 和 ROUGE-L 同样全面领先。
    \item \textbf{B2$'$v2 为最优消融配置}：Viol.\% = 4.25\%，95\% CI [3.60\%, 4.83\%]，CI 上界与其余配置 CI 下界不重叠，统计显著。
    \item \textbf{两个显著改进步骤}：A0$\to$A1（训练 $W_{\mathrm{proj}}$，$p_{\mathrm{adj}}=0.008$）和 B2$'$$\to$B2$'$v2（严格证据卡，$p_{\mathrm{adj}}=0.006$）。
    \item \textbf{R1 集中在 left/right lung}：约 28\% 的 left/right lung 句子触发 R1 侧别违规，说明路由器在肺叶层面的侧别区分仍有改进空间。
    \item \textbf{R3 集中在 bilateral}：bilateral 句子的 R3 违规率 6.0\%，因为 bilateral 解剖区域跨多个深度层级。
    \item \textbf{NLG 质量稳定}：BLEU-4 在 0.47--0.48 范围内波动极小，表明验证链不损害生成质量。
    \item \textbf{C2$'$ 出现回退}：LLM Judge 使 Viol.\% 从 4.19\% 升至 6.18\%（$p_{\mathrm{adj}}=0.006$），但 CF 从 77.6\% 升至 81.1\%，说明 Judge 提升了引用忠实度但引入了更多重路由导致的违规。
\end{enumerate}

\end{document}